\section{Certificate escape analysis}
\label{sec:escape}

The central static analysis of SCR asks what part of a certificate appears in
the public theorem.  There are three classes.

\begin{definition}[Escape classes]
Let \(c\) satisfy a contract \(V(q,c)\).
\begin{description}
\item[Sealed.]  The final formula contains no data computed from \(c\).  The
certificate is eliminated directly by \cref{thm:seal}.
\item[Observed.]  The final interface uses only an observation
\(o=O(q,c)\), and the proof establishes that all valid certificates give the
same observation, or observations within a stated diameter.
\item[Named.]  The actual code \(c\), its leastness, or a label depending on a
particular representative occurs in the output.  A direct producer or bounded
selector is mandatory.
\end{description}
\end{definition}

\subsection{The escape theorem}

\begin{theorem}[Escape analysis]
\label{thm:escape}
For a finite contract graph whose final goal is in the outer scope:
\begin{enumerate}[label=(\roman*)]
\item every sealed certificate can be removed by nested existential
      elimination, without a code bound;
\item every exactly observed certificate can be replaced by its
      single-valued output relation, without selecting a representative;
\item a named certificate can be exported only from a supplied term producer
      or from a proved bounded search.
\end{enumerate}
\end{theorem}

\begin{proof}
Part (i) is repeated application of \cref{thm:dependent-seal}.  For part (ii),
define
\[
   \mathsf{Out}(q,o)\;:\!\Longleftrightarrow\;
      \exists c\,[V(q,c)\wedge O(q,c)=o].
\]
Existence follows by opening a valid certificate and introducing its
observation.  If \(o,o'\) satisfy \(\mathsf{Out}\), open both certificates;
observer independence yields \(o=o'\).  All consumers formulated through
\(\mathsf{Out}\) therefore compose relationally.  Part (iii) is a syntactic
scope fact: an outer term cannot contain the eigenvariable introduced by
existential elimination.  To name it one must already have an outer term, or
construct one by a total bounded operation whose correctness is proved.
\end{proof}

The theorem does not assert an impossibility result for every conceivable
encoding.  It states the exact proof obligations for this natural-deduction
interface.  Any purported third method of exporting a local eigenvariable
will manifest either as an unproved choice axiom or as an omitted bound.

\subsection{Why a global master certificate is not the default}

Packing all witnesses in one natural number is sometimes valuable.  It permits
one reflected checker, one bounded search, and compact serialization.  It also
has costs:
\begin{itemize}
\item every component needs a numerical range even when it is proof-local;
\item changing one producer changes the global radix and checker;
\item a least global code imposes a canonical order unrelated to the
      mathematics;
\item relational and quotient objects are prematurely forced into named
      representatives;
\item the final proof term reflects a monolithic verifier rather than the
      logical dependency structure.
\end{itemize}
SCR therefore treats packing as a derived backend.  If every internal
certificate has a proved bound, nested scopes can be fused into one mixed-radix
code; the resulting theorem is an FQNF-style master certificate.  Conversely,
a master certificate can always be opened once and its fields used in nested
SCR scopes.  The two representations are extensionally intertranslatable when
bounds exist, but the scoped form remains meaningful when they do not.

\subsection{Sharing a certificate without leaking it}

A certificate may feed many downstream lemmas.  Open it once, prove a
conjunction of all required consequences in the same scope, and eliminate it
after the last consumer.  In Lean this amounts to constructing a local theorem
under \(V::\Gamma^{\up}\) and using \code{Prf.andI} or a downstream context.
There is no need to select the certificate merely to avoid duplicating a proof.
At the metalevel, a named Lean definition can share the derivation term.

\subsection{When observed becomes named}

Observer independence proves that a relation has at most one output, not that
there is a primitive-recursive output term.  If downstream mathematics accepts
an existentially unique output relation, the observation remains relational.
If the API insists on a term \(o(q)\), one still needs either a direct
observation producer or a bound from which an accepted certificate can be
selected.  This distinction is especially important for exact real numbers:
a unique real specified by a predicate is not automatically a
primitive-recursive Cauchy name with a modulus.
